\chapter{Conclusion and Future Work}

This chapter summarises the contributions and findings of AlgoBugs, identifies the limitations of the current study, and outlines directions for future work.

\section{Summary}

AlgoBugs is the first diagnostic benchmark designed to evaluate LLM fault-exposure capability on competitive programming bugs through structured, per-category analysis. Four primary contributions were made.

First, a novel eight-category taxonomy (T1--T8) was developed to classify competitive programming bugs by the specific reasoning capability required to expose them: T1 Integer Overflow, T2 Modular Arithmetic, T3 Off-by-One, T4 Wrong Conditional, T5 Algorithmic Flaw, T6 State/Initialization, T7 Corner Case, and T8 I/O Format. The taxonomy was validated through an inter-rater reliability study (Cohen's $\kappa$), confirming consistent applicability across annotators.

Second, a curated benchmark dataset of 1,000 paired C++ submissions was assembled from 40 Codeforces problems spanning eight difficulty brackets (rating 1400--2100). Each pair is execution-validated, manually reviewed for single-fault clarity, and labelled under the taxonomy.

Third, a reproducible evaluation engine was built to compute per-category Fault Exposure Rate (FER) for any LLM accessible via the OpenRouter API. The engine supports concurrent evaluation, partial-run recovery, and structured JSON result logging.

Fourth, five LLMs (Gemini 2.0 Flash, DeepSeek V3, LLaMA 3.3 70B, Qwen3 Coder 30B, and GPT-5.1) were evaluated under three prompting strategies: zero-shot, chain-of-thought, and few-shot, producing 15 model-strategy combinations. FER ranged from 10.8\% to 21.2\% across all combinations. T3 (Off-by-One) was the most exposed category at a 25.6\% average; T1 (Integer Overflow) was the hardest at 4.3\%. Chain-of-thought prompting showed no consistent benefit and hurt DeepSeek V3 by 3.4 percentage points. Few-shot prompting improved LLaMA 3.3 70B (+3.4 pp) and GPT-5.1 (+5.6 pp) but had no meaningful effect on the remaining three models, indicating model-specific sensitivity to task demonstrations rather than a general benefit. The consistent FER ceiling near 21\% across all configurations reflects that current LLMs generate specification-valid test inputs rather than fault-targeted ones.

\section{Limitation}

\textbf{Language restriction.} All 1,000 solution pairs in AlgoBugs are written in C++. Bug patterns in C++ are shaped by its fixed-width integer types: a 32-bit \texttt{int} overflows at approximately 2.1 billion, and T1 bugs in the dataset exploit this boundary. In Python, the built-in \texttt{int} type uses arbitrary precision, which means the entire T1 category would behave differently or cease to exist. In Java, the behavior is similar to C++ but the idioms for modular arithmetic differ. The per-category FER values reported here are specific to the C++ context; they cannot be interpreted as language-independent assessments of LLM reasoning capability.

\textbf{Model coverage.} Five models were evaluated, representing a snapshot of the LLM landscape as of early 2026. GPT-5.1 completed only 870 of the 1,000 pairs due to a project credit constraint, so its per-category results are computed over a slightly smaller denominator than the other models. More importantly, all five models are under active development; a model version released six months after this benchmark was run may produce substantially different FER values on the same dataset. The benchmark is designed to support re-evaluation as new model versions become available, but the specific numerical results in this study reflect a fixed evaluation snapshot.

\textbf{Problem difficulty range.} All problems fall within the 1400--2100 Codeforces difficulty rating, which corresponds to the upper-intermediate to advanced competitive programming range. Problems rated above 2100 were excluded because submission volume thins out at higher difficulties: fewer users attempt these problems, and the WA-to-AC pair mining strategy requires a sufficient pool of submissions per problem to find valid pairs. It is plausible that problems above 2100 involve more compositional algorithmic bugs (e.g., graph algorithms with subtle boundary conditions) that would shift the T5 and T7 category distributions. The benchmark's applicability to harder problems remains an open question.

\textbf{Few-shot demonstration design.} The three synthetic few-shot demonstrations were drawn from outside the AlgoBugs dataset and were not matched to the bug category of the test item. A category-matched design, where the demonstrations specifically show how to expose T1 bugs when the test item is also T1, was excluded to avoid a specific data-leakage risk: if a demonstration problem were semantically similar to a test problem, the model would have indirect access to the solution structure. The cost of this design choice is that the demonstrations teach the task format but not category-specific fault-exposure reasoning. This tradeoff is discussed further in the Future Work section.

\section{Future Work}

\textbf{Category-matched few-shot design.} The most direct extension of the prompting experiments is to provide demonstrations matched to the bug category of each test item. A T1 demonstration would show a solution with an integer overflow bug, a test input that produces different output on the correct and buggy versions, and an explanation of why the specific value was chosen. For T1 and T2, where average zero-shot FER is below 10\%, the model needs to see what a fault-targeted test looks like for arithmetic overflow before generating one itself. The current synthetic demonstrations anchor the output format but convey no category-specific information. Avoiding data leakage in this design requires held-out pairs not present in the 1,000-pair benchmark, which are available from the same Codeforces mining pipeline.

\textbf{Dataset expansion.} The current benchmark contains 25 pairs per bug category on average. This is sufficient for computing category-level FER but is below the threshold for reliable interaction analysis between difficulty rating and category. A dataset of 5,000 or more pairs, with at least 100 pairs per category, would enable difficulty-stratified analysis: for example, whether T3 (Off-by-One) is consistently easy across all ratings or only at lower difficulties where loop bounds are simpler. The evaluation engine supports this expansion directly; the limiting factor is manual review time for single-fault verification.

\textbf{Fine-tuning experiments.} The per-category results identify T1 (Integer Overflow) and T2 (Modular Arithmetic) as consistent weak points across all five models and all three prompting strategies. These categories represent failures of arithmetic fault-targeted reasoning rather than output format. Targeted fine-tuning on a training set of T1 and T2 pairs, followed by evaluation using the AlgoBugs engine, would directly test whether this specific capability can be improved by supervised training. Because the evaluation engine already computes per-category FER automatically, no additional instrumentation is required to measure the effect of fine-tuning.

\textbf{Public dataset and engine release.} Releasing the 1,000-pair dataset on HuggingFace and the evaluation engine on GitHub would allow the research community to evaluate new models as they are released without repeating the data collection and annotation steps. The engine's OpenRouter-based API layer makes it straightforward to add new models by updating a single configuration entry. Public release would also enable extension to Python or Java solutions using the same problem set, allowing cross-language comparison of the T1 behavior described in the limitations above. Integration with an existing evaluation leaderboard, such as the BigCode evaluation harness, would make AlgoBugs results directly comparable to other code benchmarks.

\textbf{Hint-augmented evaluation.} A separate experimental design would provide the model with the bug category before generating a test: for example, informing the model that the solution contains a T1 (Integer Overflow) bug and asking it to generate a test that exposes this specific kind of fault. Comparing fault-exposure rates in the hint condition versus the standard condition would isolate two distinct capabilities that the current task measures jointly: the ability to reason about code to identify the likely bug location, and the ability to construct a numerical test input that triggers the fault. If hint-augmented FER for T1 rises substantially above the standard 4.3\%, the bottleneck is test construction, not code reasoning. If it does not rise, the bottleneck is the reasoning step itself.
